% GENERATED FILE. Edit canonical lessons, then run npm run generate:books.
% canonical-source-hash: fnv1a64:22b07dd03e99fab6
% canonical-lessons: PT-C520-cenario, PT-C520-encenar, PT-C520-encenacao, PT-C520-estreia, PT-C520-estrear

\chapter{Estreias}
\label{ch:pt-estreias}
\hlchaptermodality{520}

\begin{chapteropening}
\textbf{By the end of this chapter:} \emph{I can talk about staging a play and opening night.}

Complete the last lesson of chapter 520: I can talk about staging a play and opening night.
\end{chapteropening}

\section[o cenário]{o cenário — set, scenery; setting}

\label{lesson:PT-C520-cenario}



\begin{quote}
Before the new one: say the Portuguese for audience, then the Portuguese for applause.
\end{quote}

\subsection*{You'll want to know: o cenário}
\textbf{o cenário} — \textquotedblleft{}set, scenery; setting\textquotedblright{}.

\begin{quote}
O cenário da peça era uma cozinha antiga.
\end{quote}

\begin{quote}
\emph{The play's set was an old kitchen.}
\end{quote}

Also a setting or a scenario: \emph{um cenário bonito}.

\begin{cousinweb}
\emph{cenário} is built on Latin \emph{scaena} \textquotedblleft{}a stage, a scene\textquotedblright{} (EtymDB), from Greek \emph{skēn\'{\={e}}} \textquotedblleft{}a tent, a stage building\textquotedblright{}.
\end{cousinweb}

\begin{grammarlens}[title={What you've built}]
One more word for telling, explaining and describing at B1.
\end{grammarlens}

\subsection*{Your turn}
\begin{itemize}
\raggedright
  \item \emph{Say it:} \emph{o cenário}
  \item \emph{Say it:} \emph{o cenário}, once more
  \item \emph{Say it:} a sentence of your own with \emph{o cenário}
  \item \emph{Recall:} say the Portuguese for audience, then the Portuguese for applause, then say \emph{o cenário} again
\end{itemize}

\begin{tcolorbox}[breakable,colback=teal!4,colframe=teal!35!black,title={Before you move on}]
What does o cenário mean? (\textquotedblleft{}Set, scenery; setting\textquotedblright{}.) Say it once more.
\end{tcolorbox}
\section[encenar]{encenar — to stage (a play)}

\label{lesson:PT-C520-encenar}



\begin{quote}
Before the new one: say the Portuguese for applause, then the Portuguese for set.
\end{quote}

\subsection*{You'll want to know: encenar}
\textbf{encenar} — \textquotedblleft{}to stage (a play)\textquotedblright{}.

\begin{quote}
O professor encenou a peça com os alunos.
\end{quote}

\begin{quote}
\emph{The teacher staged the play with the pupils.}
\end{quote}

Built on \emph{a cena} \textquotedblleft{}a scene\textquotedblright{}.

\begin{cousinweb}
\emph{encenar} is built on \emph{cena} \textquotedblleft{}a scene\textquotedblright{}, from Latin \emph{scaena} (EtymDB), which goes back to Greek \emph{skēn\'{\={e}}} \textquotedblleft{}a stage building\textquotedblright{}.
\end{cousinweb}

\begin{grammarlens}[title={What you've built}]
One more word for telling, explaining and describing at B1.
\end{grammarlens}

\subsection*{Your turn}
\begin{itemize}
\raggedright
  \item \emph{Say it:} \emph{encenar}
  \item \emph{Say it:} \emph{encenar}, once more
  \item \emph{Say it:} a sentence of your own with \emph{encenar}
  \item \emph{Recall:} say the Portuguese for applause, then the Portuguese for set, then say \emph{encenar} again
\end{itemize}

\begin{tcolorbox}[breakable,colback=teal!4,colframe=teal!35!black,title={Before you move on}]
What does encenar mean? (\textquotedblleft{}To stage (a play)\textquotedblright{}.) Say it once more.
\end{tcolorbox}
\section[a encenação]{a encenação — staging, production}

\label{lesson:PT-C520-encenacao}



\begin{quote}
Before the new one: say the Portuguese for set, then the Portuguese for to stage.
\end{quote}

\subsection*{You'll want to know: a encenação}
\textbf{a encenação} — \textquotedblleft{}staging, production\textquotedblright{}.

\begin{quote}
A encenação foi moderna e original.
\end{quote}

\begin{quote}
\emph{The staging was modern and original.}
\end{quote}

Built on \emph{encenar}.

\begin{cousinweb}
\emph{encenação} is built on \emph{encenar} \textquotedblleft{}to stage\textquotedblright{}, from \emph{cena} \textquotedblleft{}a scene\textquotedblright{}, which goes back through Latin \emph{scaena} to Greek \emph{skēn\'{\={e}}} (EtymDB).
\end{cousinweb}

\begin{grammarlens}[title={What you've built}]
One more word for telling, explaining and describing at B1.
\end{grammarlens}

\subsection*{Your turn}
\begin{itemize}
\raggedright
  \item \emph{Say it:} \emph{a encenação}
  \item \emph{Say it:} \emph{a encenação}, once more
  \item \emph{Say it:} a sentence of your own with \emph{a encenação}
  \item \emph{Recall:} say the Portuguese for set, then the Portuguese for to stage, then say \emph{a encenação} again
\end{itemize}

\begin{tcolorbox}[breakable,colback=teal!4,colframe=teal!35!black,title={Before you move on}]
What does a encenação mean? (\textquotedblleft{}Staging, production\textquotedblright{}.) Say it once more.
\end{tcolorbox}
\section[a estreia]{a estreia — premiere, opening night}

\label{lesson:PT-C520-estreia}



\begin{quote}
Before the new one: say the Portuguese for to stage, then the Portuguese for staging.
\end{quote}

\subsection*{You'll want to know: a estreia}
\textbf{a estreia} — \textquotedblleft{}premiere, opening night\textquotedblright{}.

\begin{quote}
A estreia do filme é na sexta-feira.
\end{quote}

\begin{quote}
\emph{The film's premiere is on Friday.}
\end{quote}

Also a debut: \emph{a estreia na seleção}.

\begin{cousinweb}
\emph{estreia} is inherited, through Old Portuguese \emph{estrẽa}, from Latin \emph{strēna} \textquotedblleft{}a sign, an omen\textquotedblright{}, and so a New Year gift that brought luck (EtymDB).
\end{cousinweb}

\begin{grammarlens}[title={What you've built}]
One more word for telling, explaining and describing at B1.
\end{grammarlens}

\subsection*{Your turn}
\begin{itemize}
\raggedright
  \item \emph{Say it:} \emph{a estreia}
  \item \emph{Say it:} \emph{a estreia}, once more
  \item \emph{Say it:} a sentence of your own with \emph{a estreia}
  \item \emph{Recall:} say the Portuguese for to stage, then the Portuguese for staging, then say \emph{a estreia} again
\end{itemize}

\begin{tcolorbox}[breakable,colback=teal!4,colframe=teal!35!black,title={Before you move on}]
What does a estreia mean? (\textquotedblleft{}Premiere, opening night\textquotedblright{}.) Say it once more.
\end{tcolorbox}
\section[estrear]{estrear — to premiere, to open; to use for the first time}

\label{lesson:PT-C520-estrear}



\begin{quote}
Before the new one: say the Portuguese for staging, then the Portuguese for premiere.
\end{quote}

\subsection*{You'll want to know: estrear}
\textbf{estrear} — \textquotedblleft{}to premiere, to open; to use for the first time\textquotedblright{}.

\begin{quote}
O teatro estreia uma peça nova esta semana.
\end{quote}

\begin{quote}
\emph{The theatre opens a new play this week.}
\end{quote}

Also to wear something for the first time: \emph{estreei os sapatos}.

\begin{cousinweb}
\emph{estrear} is inherited from Old Portuguese \emph{estrẽar} (EtymDB), built like \emph{estreia} on Latin \emph{strēna} \textquotedblleft{}an omen, a lucky gift\textquotedblright{}.
\end{cousinweb}

\begin{grammarlens}[title={What you've built}]
One more word for telling, explaining and describing at B1.
\end{grammarlens}

\subsection*{Your turn}
\begin{itemize}
\raggedright
  \item \emph{Say it:} \emph{estrear}
  \item \emph{Say it:} \emph{estrear}, once more
  \item \emph{Say it:} a sentence of your own with \emph{estrear}
  \item \emph{Recall:} say the Portuguese for staging, then the Portuguese for premiere, then say \emph{estrear} again
\end{itemize}

\begin{tcolorbox}[breakable,colback=teal!4,colframe=teal!35!black,title={Before you move on}]
What does estrear mean? (\textquotedblleft{}To premiere, to open; to use for the first time\textquotedblright{}.) Say it once more.
\end{tcolorbox}
